\section{Descent of preschemes affine over another one}
\label{VIII.2}
% original p. 202

Since the inverse image functor for Modules is compatible with the
tensor product and with other tensor operations,
Theorem~\Ref{VIII.1.1} implies various variants, obtained by
considering, instead of a single quasi-coherent Module, a
quasi-coherent Module or a system of quasi-coherent Modules endowed
with various additional structures expressible by means of tensor
operations. For example, the datum of three quasi-coherent
Modules~$F$, $G$, $H$ on~$S$ and of a pairing
$$
F\otimes G\to H\quoi,
$$
is equivalent to the datum of three quasi-coherent Modules~$F'$,
$G'$, $H'$ on~$S'$, endowed with descent data relative
to~$g\colon S'\to S$, and endowed with a pairing
$$
F'\otimes G'\to H'\quoi,
$$
``compatible'' with these descent data, in the obvious sense of the
term. For example, if $F=G=H$, one sees that the datum of a
quasi-coherent Module~$F$ on~$S$ endowed with an algebra law (which
for the moment we do not assume to satisfy any axiom of
associativity, of commutativity or of existence of a unit section),
is equivalent to the same datum on~$S'$, endowed in addition with a
descent datum. And using the results of the preceding no., one
finds at once that for~$F$ to satisfy one of the usual axioms
alluded to just now, it is necessary and sufficient that the same
hold for~$F'$. For example, the datum of a quasi-coherent
Algebra~$\mathcal{A}$ on~$S$ (by which we henceforth understand:
associative, commutative, with unit section) is equivalent to the
datum of a quasi-coherent Algebra~$\mathcal{A}'$ on~$S'$, endowed
with a descent datum relative to
$g\colon S'\to S$.
If one recalls the equivalence between the dual category of the
quasi-coherent Algebras on~$S$, and the category of
the~$S$-preschemes affine over~$S$, (EGA~II par.1), one at once
finds:

\begin{theorem}
\label{VIII.2.1}
Let $\cal{F}'$ be the fibered category of affine morphisms of
preschemes $f\colon X\to S$, considered as a fibered subcategory of
the fibered category
% original p. 203
of arrows in the category of preschemes $\Sch$
\textup{(VI~\Ref{VI.11}.a)}. Let $g\colon S'\to S$ be a faithfully
flat and quasi-compact morphism of preschemes. Then $g$ is an
effective $\cal{F}'$-descent morphism.
\end{theorem}

\section[Descent of set-theoretic properties]{Descent of set-theoretic properties and of finiteness properties of morphisms\protect\footnotemark}
\label{VIII.3}

\footnotetext{For other results like those treated in
nos.~3 and~4, \Cf EGA~IV~2.3, 2.6,~2.7.}

\begin{proposition}
\label{VIII.3.1}
Let $f\colon X\to Y$ be an~$S$-morphism, $g\colon S'\to S$ a
surjective morphism, $f'\colon X'=X\times_{S}S'\to Y'=Y\times_{S}S'$
the morphism deduced from~$f$ by the change of base by means
of~$g\colon S'\to S$. For $f$ to be surjective (\resp radicial), it
is necessary and sufficient that $f'$ be so.
\end{proposition}

We note that $f'$ can also be obtained by the change of base~$Y'\to
Y$, which is likewise surjective since it is deduced from~$g\colon
S'\to S$, which is. On the other hand, for every~$y\in Y$ and
every~$y'\in Y'$ over~$y$, we have an isomorphism
$$
X'_{y'}\isomfrom X_{y}\otimes_{\kres(y)}\kres(y')\quoi,
$$
where $X_{y}$ denotes the fiber of~$X$ at~$y$, and $X'_{y'}$ that
of~$X'$ at~$y'$. It follows that~$X_{y}$ is nonempty (\resp has at
most one point, and the latter corresponds to a radicial residue
field extension) if and only if
$X'_{y'}$
has the same property. This proves~\Ref{VIII.3.1}.

\begin{corollary}
\label{VIII.3.2}
Under the conditions of~\Ref{VIII.3.1}, if $f'$ is injective
(\resp bijective), then $f$ is so as well.
\end{corollary}

This comes from the fact that if $X'_{y'}$ has at most one point
(\resp exactly one point) the same holds for~$X_{y}$; this is indeed
so, since the morphism $X'_{y'}\to X_{y}$ is surjective (for it is
deduced from~$\Spec(\kres(y'))\to \Spec(\kres(y))$, which is).

\begin{proposition}
\label{VIII.3.3}
With the notation of~\Ref{VIII.3.1}. Suppose that~$g\colon S'\to S$
is surjective and quasi-compact (\resp faithfully flat and
quasi-compact). For~$f$ to be quasi-compact (\resp of finite type),
it is necessary and sufficient that~$f'$ be so.
\end{proposition}

Only
% original p. 204
the~``it suffices'' needs to be proved. One may evidently
assume~$S=Y$, since the hypothesis made on~$g\colon S'\to S$ is
preserved for~$Y'\to Y$. Moreover one may assume~$Y$ affine. Then
$Y'$ is quasi-compact, hence $X'$ is quasi-compact (since~$f'$ is so
by hypothesis). Let $(X_{i})_{i\in I}$ be a family of affine open
sets of~$X$ covering~$X$; then the~$X_{i}'$ are open sets of~$X'$
covering~$X'$, hence there is a finite subfamily which covers~$X'$.
As $X'\to X$ is surjective, it follows that the corresponding~$X_{i}$
already cover~$X$, hence~$X$ is quasi-compact, \ie $f$ is
quasi-compact. Suppose now $f'$ of finite type, and let us prove
that~$f$ is so, assuming~$g$ faithfully flat. Replacing $Y'$ by the
sum of a family of affine open sets covering it, one may assume $Y'$
affine. Finally, $X$ being covered by a finite number of affine open
sets $X_{i}$ by what precedes, one must show that they are of finite
type over~$Y$, knowing that $X_{i}'$ is of finite type over~$Y'$.
This then brings us back to the

\begin{corollary}
\label{VIII.3.4}
Let $B$ be an $A$-algebra, $A'$ a faithfully flat~$A$-algebra,
$B'=B\otimes_{A}A'$ the~$A$-algebra deduced from~$B$ by change of
ring. For~$B$ to be of finite type, it is necessary and sufficient
that~$B'$ be so.
\end{corollary}

Only the~``it suffices'' needs to be proved. We have
$B=\varinjlim_i B_{i}$, where the $B_{i}$ are the subalgebras of
finite type of~$B$. Hence we have
$B'=\varinjlim_i B_{i}'$, and if $B'$ is of finite type
over~$A'$, then $B'$ is equal to one of the~$B_{i}'$, hence $B$ is
equal to~$B_{i}$, hence is of finite type.

\begin{corollary}
\label{VIII.3.5}
Suppose again that the change of base morphism~$g\colon S'\to S$ is
faithfully flat and quasi-compact. For~$f$ to be quasi-finite, it is
necessary and sufficient that~$f'$ be so.
\end{corollary}

Indeed, the property ``quasi-finite'' is by definition the
conjunction of ``of finite type'' and ``with finite fibers'', each
of which descends well by~$g$, the first by virtue
of~\Ref{VIII.3.3}, the second by the reasoning of~\Ref{VIII.3.1}
(which uses only the fact that $g$ is surjective).

\begin{remarks}
\label{VIII.3.6}
Let $A$ be a ring, $X$ an~$A$-prescheme. One easily sees that
% original p. 205
the following conditions are equivalent:

\begin{enumerate}
\item[(i)] There exist a noetherian ring~$A_{0}$ (which one may, if
one wishes, assume to be a subring of finite type of~$A$), an
$A_{0}$-prescheme of finite type~$X_{0}$, a homomorphism
$A_0\to A$
and an $A$-isomorphism~$X\isomto X_{0}\times_{A_{0}}A$.
\item[(ii)] The diagonal morphism~$X\to X\times_{\Spec(A)}X$ is
quasi-compact (an empty condition if $X$ is separated over~$A$), $X$
is a finite union of affine open sets~$X_{i}$ whose rings~$B_{i}$
are algebras of finite presentation over~$A$, \ie quotients of
polynomial algebras in a finite number of indeterminates, by ideals
of finite type.
\end{enumerate}

If $X$ is itself affine, with ring~$B$, these conditions simply
mean that~$B$ is an algebra of finite presentation over~$A$.

A morphism~$f\colon X\to Y$ is called a \emph{morphism of finite
presentation}, and one also says that~$X$ is of finite presentation
over~$Y$, if $Y$ is a union of affine open sets~$Y_{i}$ such
that~$X|Y_{i}$, as a $Y_{i}$-prescheme, satisfies the preceding
equivalent conditions. The same then holds for $X|Y'$ for
\emph{every} affine open set $Y'$ in $Y$. This is a property stable
under change of base, and moreover the composite of two morphisms of
finite presentation is of finite presentation.

These notions having been set up, one sees on (ii), proceeding as
in~\Ref{VIII.1.10}, that this statement remains valid when the words
``of finite type'' are replaced in it by ``of finite presentation''.
\end{remarks}

\section{Descent of topological properties}
\label{VIII.4}

\begin{theorem}
\label{VIII.4.1}
Let $g\colon Y'\to Y$ be a morphism, and $Z$ a subset of~$Y$. Assume
that~$g$ is flat, and that there exists a quasi-compact
morphism~$f\colon X\to Y$ such that~$Z=f(X)$ (N.B. if $Y$ is
noetherian, this last condition is implied by
``$Z$ is constructible''). Then we have
$$
g^{-1}(\overline{Z})=\overline{g^{-1}(Z)}\quoi.
$$
\end{theorem}

One may assume~$Y$ affine, then $Y'$ affine. As $Y$ is affine,
$X$ is a finite
% original p. 206
union of affine open sets~$X_{i}$, and replacing~$X$ by the sum
of the~$X_{i}$, one may also assume~$X$
affine. Let~$A,A',B$ be the rings of~$Y,Y',X$, $B'=B\otimes_{A}A'$
that of~$X'=X\times_{Y}Y'$, $I$ the kernel of~$A\to B$, $I'$ the
kernel of~$A'\to B'$; thus the closed subsets of~$Y$ and~$Y'$
defined by these ideals are respectively the closure
of~$Z=f(X)$ and the closure of~$Z'=f'(X')=g^{-1}(Z)$. We want to
establish that the latter is equal to
$g^{-1}(\overline{Z})$, which will follow from~$I'=IA'$, itself a
consequence of the flatness of~$A'$ over~$A$.

\begin{corollary}
\label{VIII.4.2}
Let $g\colon Y'\to Y$ be a flat and quasi-compact morphism, and $Z'$
a closed subset of~$Y'$ saturated for the set-theoretic equivalence
relation defined by $g$. Then we have
$Z'=g^{-1}(\overline{g(Z')})$.
\end{corollary}

Indeed we have $Z'=g^{-1}(Z)$, with~$Z=g(Z')$. One can then
apply~\Ref{VIII.4.1}, noting that the condition imposed on~$Z$
in~\Ref{VIII.4.1} is indeed satisfied by taking for~$X$ the
prescheme~$Z'$ endowed with the reduced structure induced
by~$Y'$. (The fact that~$g$ is quasi-compact then ensures that~the
induced morphism~$f\colon Z'\to Y$ is quasi-compact).

Statement~\Ref{VIII.4.2} also means that \emph{the topology
of~$g(Y')$ induced by~$Y$ is the quotient of that of~$Y'$}. In
particular:

\begin{corollary}
\label{VIII.4.3}
Let $g\colon Y'\to Y$ be a faithfully flat and quasi-compact
morphism. Then $g$ makes~$Y$ a quotient topological space
of~$Y'$, \ie for a subset~$Z$ of~$Y$, $Z$ is closed
(\resp open) if and only if $Z'=g^{-1}(Z)$ is.
\end{corollary}

Let us now recall that two elements $a,b$ of~$Y'$ have the same
image in $Y$ if and only if they are of the form
$p_{1}(c),p_{2}(c)$ for a suitable element $c$
of~$Y''=Y'\times_{Y}Y'$. It follows that if $g$ is surjective,
we have an \emph{exact} diagram of sets
$$
\xymatrix@C=.5cm{\cal{P}(Y)\ar[r] & \cal{P}(Y')\ar@<2pt>[r]\ar@<-2pt>[r] &
\cal{P}(Y'')}\quoi,
$$
where for every set~$E$, we denote by~$\cal{P}(E)$ the set of its
subsets. This being so, \Ref{VIII.4.3} can also be interpreted as
follows:

\begin{corollary}[Descent of open \resp closed subsets]
\label{VIII.4.4}
Let
% original p. 207
$g\colon Y'{\to} Y$ be as in~\Ref{VIII.4.3}. For every
prescheme~$X$, let $\Ouv(X)$ \resp $\Fer(X)$ be the set of its open
subsets \resp the set of its closed subsets. Then we have
\emph{exact} diagrams of maps of sets (deduced from~$g$ and from the
two projections of $Y''=Y'\times_{Y}Y'$):
\vspace*{-5pt}
\begin{gather*}
\xymatrix@C=.5cm{\Ouv(Y)\ar[r] & \Ouv(Y')\ar@<2pt>[r]\ar@<-2pt>[r] &
\Ouv(Y'')}\\[-3pt]
\xymatrix@C=.5cm{\Fer(Y)\ar[r] & \Fer(Y')\ar@<2pt>[r]\ar@<-2pt>[r] &
\Fer(Y'')}
\end{gather*}
\end{corollary}

We have the following complement to~\Ref{VIII.4.3}:

\begin{corollary}
\label{VIII.4.5}
Let $g\colon Y'\to Y$ be as in~\Ref{VIII.4.3}, and let~$Z$ be a
subset of~$Y$ such that there exists a quasi-compact
morphism~$f\colon X\to Y$ with image~$Z$ (for example $Z$
constructible, $Y$ noetherian). For $Z$ to be a locally closed
subset of~$Y$, it is necessary and sufficient that $Z'=g^{-1}(Z)$ be
a locally closed subset of~$Y'$.
\end{corollary}

It suffices to prove the ``it suffices''. Let $Y_{1}$ be the closed
subprescheme of~$Y$, closure of~$Z$ endowed with the induced reduced
structure, and let~$Y_{1}'=Y_{1}\times_{Y}Y'$ be the closed
subprescheme of~$Y'$ inverse image of~$Y_{1}$. Its underlying set is
the inverse image~$g^{-1}(Y_{1})=g^{-1}(\overline{Z})$, hence is
equal, by virtue of~\Ref{VIII.4.1}, to~$\overline{Z'}$. As $Z'$ is
locally closed in~$Y'$, it is open in~$\overline{Z'}$, hence open
in~$Y_{1}'$. Now the latter is faithfully flat and quasi-compact
over~$Y_{1}$, hence by virtue of~\Ref{VIII.4.3} one concludes
that~$Z$ is open in~$Y_{1}$, \ie in~$\overline{Z}$, which means that
$Z$ is locally closed.

\begin{corollary}
\label{VIII.4.6}
Let $g\colon S'\to S$ be a faithfully flat and quasi-compact
morphism, $f\colon X\to Y$ an $S$-morphism, $f'\colon X'\to Y'$ the
$S'$-morphism deduced from it by change of base. Suppose that~$f'$ is
an open map (\resp a closed map, \resp quasi-compact and a
homeomorphism into, \resp a homeomorphism onto); then $f$ has the
same property.
\end{corollary}

As $Y'$ is faithfully flat and quasi-compact over~$Y$, one may
assume~$Y=S$. Let~$Q$
% original p. 208
be a subset of~$X$; we then have (denoting by $h$ the projection
morphism $X'\to X$):
$$
g^{-1}(f(Q))=f'(h^{-1}(Q))\quoi.
$$

If $Q$ is open (\resp closed), the same holds for~$h^{-1}(Q)$, hence
also for~$f'(h^{-1}(Q))$ if $f'$ is assumed to be an open map
(\resp a closed map), hence the same holds for~$f(Q)$ by virtue of the
preceding formula and of~\Ref{VIII.4.3}. This proves the first two
assertions in~\Ref{VIII.4.6}; it remains to examine the case
where~$f'$ is a homeomorphism into, and to prove then that~$f$ is a
homeomorphism into. (The case of a homeomorphism onto will then
follow from~\Ref{VIII.3.1}). By virtue of~\Ref{VIII.3.2}
$f$ is injective; it remains to prove that the map $X\to f(X)$ is
open. We already know that~$f$ is quasi-compact by virtue
of~\Ref{VIII.3.3}. It suffices to prove that for every closed
subset~$Z$ of~$X$, we have $Z=f^{-1}(\overline{f(Z)})$, which is
equivalent to the analogous formula for the inverse images under the
surjective map~$h\colon X'\to X$, \ie to
$$
Z'=f'^{-1}(g^{-1}(\overline{f(Z)})\quoi,
$$
where we set $Z'=h^{-1}(Z)$. Now by virtue of~\Ref{VIII.4.1} applied
to the subset~$f(Z)$ of~$Y$, we have
$g^{-1}(\overline{f(Z)})=\overline{g^{-1}(f(Z))}$, and the formula to
be proved is equivalent to
$$
Z'=f'^{-1}(\overline{f'(Z')})\quoi,
$$
which follows from the hypothesis that $f'$ is a homeomorphism
into.

N.B. In this last argument, having already assumed that~$f$ is
quasi-compact, we did not use that~$g$ is quasi-compact, but only
that
$g$
is faithfully flat. Hence it is under this hypothesis that one can
descend the property ``homeomorphism into'', or ``homeomorphism
onto'', or also, thanks to the preceding argument, the property
``$f'$ is quasi-compact and makes~$f'(X')$ a quotient topological
space of~$X'$''.

We shall say that a morphism~$f\colon X\to Y$ of preschemes is
\emph{universally open} (\resp \emph{universally closed},
\resp \emph{universally bicontinuous}, etc.)
if for every change of base~$Y'\to Y$, $f'\colon X'\to Y'$ is an open
morphism (\resp closed, \resp a
homeomorphism
onto the image space). One then derives from~\Ref{VIII.4.6}:

\begin{corollary}
\label{VIII.4.7}
Under
% original p. 209
the conditions of~\Ref{VIII.4.6}, for~$f$ to be universally open,
(\resp universally closed, \resp a universal homeomorphism into,
\resp a universal homeomorphism), it is necessary and sufficient
that $f'$ be so.
\end{corollary}

\begin{corollary}
\label{VIII.4.8}
Under the conditions of~\Ref{VIII.4.6}, for~$f$ to be separated
(\resp proper) it is necessary and sufficient that $f'$ be so.
\end{corollary}

To say that~$f$ is separated means that the diagonal
morphism~$X\to X\times_{Y}X$ is closed, or also universally closed,
and the first assertion
of
\Ref{VIII.4.8} therefore follows from~\Ref{VIII.4.7}. To say
that~$f$ is proper means that $f$ satisfies the conditions a) $f$ is
of finite type b) $f$ is separated c) $f$ is universally closed.
Condition a) descends well by virtue of~\Ref{VIII.3.3}, b) also
according to what we have just seen, and finally c) likewise
by~\Ref{VIII.4.7}.

\begin{remarks}
\label{VIII.4.9}
Let us recall that when $g\colon Y'\to Y$ is a flat morphism of
finite type, with~$Y$ locally noetherian, then $g$ is an open
morphism (VI~\Ref{IV.6.6}), which is a more precise result
than~\Ref{VIII.4.3}. One will note however that if~$f$ is a
faithfully flat and quasi-compact morphism of noetherian preschemes,
then $f$ is not in general an open morphism. Let for example $Y$ be
an irreducible scheme whose generic point~$y$ is not open (for
example an algebraic curve), and take for~$Y'$ the sum
scheme~$Y\amalg \Spec(\kres(y))$; then the image under the
structural morphism~$Y'\to Y$ of the open subset~$\Spec(\kres(y))$
is not an open subset of~$Y$. The reader will also notice that
various statements of the present expos\'e become false if one drops
from them the hypothesis that the faithfully flat morphism
considered is also quasi-compact, the typical case in which the
statements fail being that where one takes for~$Y'$ the sum scheme
of the spectra of the local rings of the points of~$Y$. For example,
taking again for~$Y$ an irreducible algebraic curve, and for~$Z$ the
subset of~$Y$ reduced to the generic point, its inverse image
in~$Y'$ is open, without~$Z$ being open.
\end{remarks}

\subsection{}
\label{VIII.4.10}
Various
% original p. 210
statements of the present expos\'e remain valid when the hypothesis
that~$Y'$ is flat over~$Y$ is replaced in them by the following:
there exists a Module of finite type~$F$ on~$Y'$, with support~$Y'$,
flat with respect to~$Y$; the hypothesis of faithful flatness will
then be replaced by the preceding one, plus the hypothesis that
$Y'\to Y$ is surjective. This applies to the first two assertions
in~\Ref{VIII.1.10}, to~\Ref{VIII.3.3}, \Ref{VIII.3.5},
\Ref{VIII.4.1}, and consequently to all the results of the
present no.
